\section*{四、实验心得与思考题}

\subsection*{1. 实验心得}

通过本次实验，我学会了用替代法和中值法测量电流计内阻，并利用分流、分压电阻将表头改装为电流表和电压表。在电路设计时，需要分清原表头电流、改装表总电流、标准表电流，避免参数混用。从校验结果中，我意识到即使满偏点调准后，中间刻度仍可能存在偏差，因此应检查多个校验点，再评定仪表等级。组装欧姆表时，应先短接调零，并结合表头电流半偏点理解电阻刻度的分布。此外，记录电阻箱读数时要写清小数位。读所有表时都要正视刻度，减小读数误差。

\subsection*{2. 思考题}

\begin{enumerate}
\item \textbf{为什么不能用万用表欧姆挡测量电源内阻？} 欧姆挡依靠内置电源产生测量电流。若接入另一个电源，其电动势也会影响电流，使读数失去意义，还可能损坏表头。因此，电源内阻应通过负载变化时的端电压和电流来测量。

\item \textbf{数字万用表与指针式万用表有什么差异？} 指针式万用表通过指针偏转显示数值，便于观察变化趋势，但需要估读，存在视差，欧姆挡刻度也不均匀。数字万用表将测量信号转换为数字直接显示，读数方便，电压挡通常具有较高输入电阻，对被测电路的影响较小；不过，显示位数多并不代表准确度一定高，仍需查看具体指标。

\item \textbf{欧姆表中干电池内阻变化有什么影响？} 电池内阻增大时，测量电流减小，指针偏转变小，按原刻度读出的电阻会偏大。若电动势不变且调节范围足够，可通过重新调零补偿；若电池老化同时使电动势下降，即使能调零，刻度关系也可能改变，应及时更换电池。
\end{enumerate}
